\subsection{Sensitivity Analysis}
\label{subsec:sensitivity}

To comprehensively evaluate the robustness of \textsc{ECPFL}, we conduct a sensitivity analysis on the Fashion-MNIST dataset.
We systematically investigate the impact of voting mechanisms, client selection strategies, local sparsity ratios, and prefetch window sizes on model convergence.

\begin{figure}[t]
  \centering
  % 第一行：(a) (b)
  \begin{minipage}{0.48\linewidth}
    \centering
    \includegraphics[width=\linewidth]{../figures/fashion_mnist_vote_accuracy_grouped.pdf}
    \subcaption{AC under different voting mechanisms.}
    \label{fig:vote}
  \end{minipage}
  \hfill
  \begin{minipage}{0.48\linewidth}
    \centering
    \includegraphics[width=\linewidth]{../figures/fashion_mnist_client_selection_accuracy_nature.pdf}
    \subcaption{AC under different client selection strategies.}
    \label{fig:select}
  \end{minipage}

  \vspace{1em} % 行间距

  % 第二行：(c) (d)
  \begin{minipage}{0.48\linewidth}
    \centering
    \includegraphics[width=\linewidth]{../figures/fashion_mnist-cnn_fmnist_accuracy_curves_vibrant_marker.pdf}
    \subcaption{AC under different sparsity ratios.}
    \label{fig:sparsity}
  \end{minipage}
    \hfill
     \begin{minipage}{0.48\linewidth}
    \centering
    \includegraphics[width=\linewidth]{../figures/accuracy_curves_by_window.pdf}
    \subcaption{AC under different fetch window sizes.}
    \label{fig:window}
  \end{minipage}
  \caption{Comprehensive evaluation of system configurations on Fashion-MNIST.}
  \label{fig:all_results}
\end{figure}

\subsubsection{Impact of Voting Mechanisms}

Fig.~\ref{fig:all_results}(a) compares voting mechanisms.
Union and rank-based variants (RK-0.3 to RK-0.9) converge rapidly to approximately 90\% within 20 rounds, with final accuracy above 89\%.
Threshold voting with \textsc{TH-0.7} and \textsc{TH-0.9} degrades progressively: \textsc{TH-0.9} plateaus near 84\%, a 6\% accuracy gap compared with Union.
This gap arises because a 90\% consensus threshold is overly rigid and discards class-specific gradient features.
\textsc{ECPFL-RK} avoids this bottleneck by retaining globally significant updates without inflexible frequency cutoffs.

\subsubsection{Impact of Client Selection Strategies}

Fig.~\ref{fig:all_results}(b) evaluates seven representative client selection strategies: random, sticky, sketch\_clustering, poc, tifl, haccs, and oort.
All trajectories overlap almost exactly and converge to approximately 86\%, with a final accuracy deviation below 1\%.
This confirms that \textsc{ECPFL} is insensitive to selection bias: the ECP module allows infrequently selected or low-bandwidth clients to prefetch global ciphertexts during idle periods and fetch only minimal residuals when selected.
As a result, resource-aware selection can accelerate wall-clock training without sacrificing model accuracy.

\subsubsection{Impact of Local Sparsity Ratio}

Fig.~\ref{fig:all_results}(c) shows the impact of local sparsity ratios.
Retention ratios of 0.3 and above closely match the dense baseline (sparsity=1) and converge to approximately 90\%.
At sparsity=0.1, however, accuracy plateaus near 80\%, a 10\% drop, because excessive pruning discards too much gradient information before global consensus.
A retention ratio of at least 0.3 is therefore necessary to filter local noise while preserving enough information for voting.

\subsubsection{Impact of Prefetch Window Size $R$}

Fig.~\ref{fig:all_results}(d) compares prefetch window sizes.
For $R \leq 7$, the accuracy curves are nearly identical, with peak deviation under 0.3\%, because asynchronous prefetching decouples transmission from computation and prevents critical staleness.
At $R=9$, a noticeable degradation emerges, exposing the upper bound beyond which excessive stale gradient accumulation hurts convergence.
